% Pursuing Stacks (cotes 134-2 to 134-8): archive cote 134-5, site page 25 (PDF page 26), typescript page 212.
% typescript. Draft compared: ps4.tex:1132-1189.
% AI-transcribed by Claude Opus 5.5 (claude-opus-5-5, adaptive thinking, effort high), Message Batches API, batch ps2 (msgbatch_01NVwr1Yuo8qJf5rykeGBXPK), 2026-11-12; not yet checked by eye.
\scanpage{134-5/25}
\mspage{212}
to look up in this respect,  is surely D(T),  defined say à la Illusie.

In the present context,  the main point of the property of \uline{local asphe-}%
\uline{ricity} for a topos (cf.  section 35)  is that the corresponding prohomotopy
type is essentially constant,  i.e. the topos defines an actual homotopy
type.  Thus locally aspheric topoi  and weak equivalences between these
should be eligible models for homotopy types  (more accurately,  make up
a modelizer)\hand{,} just as the basic modelizer (Cat) contained in it.  The cor-%
responding statements,  when introducing a basic localizer  $\underline{W}$ ,  should be
equally valid.  One might expect too a relative variant for the notion of
local asphericity or $\underline{W}$-asphericity,  in case of a topos X over a given
base topos T (whom we may have to suppose already locally aspheric),  with
the \deleted{property}\inserted{implication} that the corresponding object of Pro(D(T)) should be again
essentially constant.

A last question I would like to mention here is about the meaning of
the notion of a so-called ``modelizing topos'',  introduced in a somewhat
formal way in section 35,  as a \unsure{loc ally|locally} aspheric \inserted{and aspheric} topos T such that the
Lawvere element $L_T$ of T is aspheric over the final object. (We assumed at
first,  moreover,  that T be even totally\xout{x} aspheric,  but soon after the
point of view \inserted{and terminology} shifted a little and the totally aspheric case was referred \xout{t}
to as a \uline{strictly} modelizing topos\hand{,} \inserted{cf. page 68.})  As made clear there,  the expectation
suggested by the terminology is of course that such a topos should indeed
be a modelizer,  when endowed with the usual notion of weak equivalence.
This still makes sense and seems plausible enough,  when usual weak equi-%
valence is replaced by the weaker notion defined in terms of an arbitrary
basic localizer $\underline{W}$ .   A related question is to get a feeling for how res-%
trictive the conditions put on a modelizing topos \xout{\illegible{1}} are.  It is clear
that nearly all topoi met with in practice,  including those associated to
the more common topological spaces (such as locally contractible ones) are
locally aspheric  - \inserted{but} what about the condition on the Lawvere element ?
For instance, taking a topological space admitting a \inserted{finite} triangulation and
which is aspheric i.e. contractible,  is the corresponding topos \xout{modeli}
modelizing ?

\margin{25}
